% Claim statement file for row claim_3_25 -- "ISigmaSeparation".
% Auto-generated stub by scaffold/solve_chapter.py at row start. The
% manager agent (or a worker it dispatches) fills the body below with the
% claim's statement(s) from the lecture notes -- statement only, no proof
% (proofs live in the sibling `claim_3_25_proof_*.tex` file).
%
% This file is a sibling of `main.tex` inside `<subsection>/tex/`. The
% `subfiles` package lets it render both standalone and as part of
% `main.tex`. Note: `main.tex` does NOT \subfile this statement file -- the
% sibling `_proof_` file restates the statement above the proof, so
% including both in `main.tex` would be redundant. This file exists for
% standalone reading / grep convenience.

\documentclass[main]{subfiles}

\begin{document}
% ref: claim_3_25 (statement)
% title: ISigmaSeparation
\def\rowref{claim\_3\_25}
\phantomsection\label{claim_3_25}

\begin{Lem}[$i\sigma$-separation under marginalization]\label{lem:stability_separation_marginalization}
    Let $G = (J, V, E, L)$ be a CDMG (def \ref{def-cdmg}); throughout we use the notation of def \ref{not-cdmg}, the walk definition of def \ref{def:walks}, item~i., the family-relationships definition of def \ref{def_3_5} (we write $\Anc^H(C) := \bigcup_{c \in C} \Anc^H(c)$ for the ancestor set of a subset $C$ in an arbitrary CDMG $H$, using def \ref{def_3_5}, item~iv., per the convention recorded in the ``Ancestor set of $C$'' paragraph of def \ref{def:sigma_blocking}), the collider/non-collider classification of positions on a walk of def \ref{def:collider_noncollider}, item~ii., the blockable/unblockable non-collider refinement of def \ref{def:unblockable_noncollider}, the $\sigma$-blocking predicate ``$\pi$ is $C$-$\sigma$-blocked'' (equivalently ``$\sigma$-blocked given $C$'') of def \ref{def:sigma_blocking} (``$C$-$\sigma$-blocked walk'' paragraph), the marginalization operator $G^{\sm \cdot}$ of def \ref{def:G_marginalization}, and the $i\sigma$-separation predicate of def \ref{def:sigma_separation}, item~1. Let
    \[
        A,\, B,\, C \;\ins\; J \cup V
        \qquad\text{and}\qquad
        D \;\ins\; V
    \]
    be subsets of nodes. The triple $A$, $B$, $C$ is \emph{not} assumed to be pairwise disjoint, nor disjoint from the input-node set $J$, nor non-empty (matching the permissiveness of def \ref{def:sigma_separation}, item~1, on its three argument-sets). The typing $D \ins V$ is the precondition of the marginalization operator of def \ref{def:G_marginalization} (which is only defined when its target is a subset of the output-node set $V$); the case $D = \emptyset$ is admitted (then $G^{\sm D} = G$ verbatim under def \ref{def:G_marginalization}, items~i.--iv., and the equivalence below trivialises). Suppose, in addition, that
    \[
        D \cap \lp A \cup B \cup C \rp \;=\; \emptyset.
    \]

    \emph{Recall: the marginalized CDMG $G^{\sm D}$ (def \ref{def:G_marginalization} with $W = D$).}
    The marginalization $G^{\sm D}$ is the CDMG
    \[
        G^{\sm D} \;=\; \lp J^{\sm D},\, V^{\sm D},\, E^{\sm D},\, L^{\sm D} \rp,
    \]
    whose four components are given by items~i.--iv.\ of def \ref{def:G_marginalization}:
    \[
        J^{\sm D} \;=\; J,
        \qquad
        V^{\sm D} \;=\; V \sm D,
    \]
    and $E^{\sm D} \ins (J \cup (V \sm D)) \times (V \sm D)$, $L^{\sm D} \ins (V \sm D) \times (V \sm D)$ are the directed- and bidirected-edge sets witnessed, respectively, by the directed-walk-through-$D$ predicate $\Phi_E$ of def \ref{def:G_marginalization}, item~iii., and the bifurcation-through-$D$ predicate $\Phi_L$ of def \ref{def:G_marginalization}, item~iv.\ (with the standing $\ul{v} \ne \ol{v}$ constraint on $L^{\sm D}$ and the end-node arrowhead constraint of clause~(f) of item~iv.). The carrier of $G^{\sm D}$ is
    \[
        J^{\sm D} \cup V^{\sm D} \;=\; J \cup (V \sm D),
    \]
    and the input-node set is preserved verbatim: $J^{\sm D} = J$.

    \emph{Well-typedness of $A$, $B$, $C$ as subsets of the carrier of $G^{\sm D}$.}
    From $A, B, C \ins J \cup V$, the disjointness hypothesis $D \cap (A \cup B \cup C) = \emptyset$, $D \ins V$, and $J \cap V = \emptyset$ (def \ref{def-cdmg}) it follows that
    \[
        A,\, B,\, C \;\ins\; (J \cup V) \sm D \;=\; J \cup (V \sm D) \;=\; J^{\sm D} \cup V^{\sm D}.
    \]
    Hence the $i\sigma$-separation predicate of def \ref{def:sigma_separation}, item~1, instantiated on the CDMG $G^{\sm D}$ with the same triple $(A, B, C)$ is well-typed on the right-hand side of the equivalence below (its argument-sets are subsets of the carrier of the CDMG on which it is evaluated).

    \emph{The equivalence (with $i\sigma$-separation unfolded inline on both sides per def \ref{def:sigma_separation}, item~1).}
    The following two universal statements are equivalent:
    \begin{enumerate}
        \item For every integer $n \ge 0$ and every walk
            \[
                \pi = (v_0, a_0, v_1, a_1, \dots, v_{n-1}, a_{n-1}, v_n)
            \]
            in $G$ in the sense of def \ref{def:walks}, item~i.\ (so $v_0, \dots, v_n \in J \cup V$ and $a_0, \dots, a_{n-1} \in E \cup L$, satisfying the walk constraint at every $k \in \{0, 1, \dots, n - 1\}$ relative to $E, L$:
            \[
                \bigl(a_k = (v_k, v_{k+1}) \;\land\; a_k \in E \cup L\bigr)
                \;\lor\;
                \bigl(a_k = (v_{k+1}, v_k) \;\land\; a_k \in E\bigr)\,)
            \]
            with $v_0 \in A$ and $v_n \in J \cup B$, the walk $\pi$ is $C$-$\sigma$-blocked in $G$ in the sense of def \ref{def:sigma_blocking} (``$C$-$\sigma$-blocked walk'' paragraph): there exists a position $k \in \{0, 1, \dots, n\}$ such that either ($k$ is a collider on $\pi$ in the sense of def \ref{def:collider_noncollider}, item~ii.\ computed relative to the edges of $G$, and $v_k \notin \Anc^G(C)$), or ($k$ is a blockable non-collider on $\pi$ in the sense of def \ref{def:unblockable_noncollider} computed relative to the edges of $G$, and $v_k \in C$); the ancestor set on the right-hand side of each existential clause is $\Anc^G(C) = \bigcup_{c \in C} \Anc^G(c)$ from def \ref{def_3_5}, item~iv., computed in $G$.

        \item For every integer $n \ge 0$ and every walk
            \[
                \pi = (v_0, a_0, v_1, a_1, \dots, v_{n-1}, a_{n-1}, v_n)
            \]
            in $G^{\sm D}$ in the sense of def \ref{def:walks}, item~i.\ \emph{applied to the CDMG $G^{\sm D}$ in place of $G$} (so $v_0, \dots, v_n \in J^{\sm D} \cup V^{\sm D} = J \cup (V \sm D)$ and $a_0, \dots, a_{n-1} \in E^{\sm D} \cup L^{\sm D}$, satisfying the walk constraint at every $k \in \{0, 1, \dots, n - 1\}$ relative to $E^{\sm D}, L^{\sm D}$:
            \[
                \bigl(a_k = (v_k, v_{k+1}) \;\land\; a_k \in E^{\sm D} \cup L^{\sm D}\bigr)
                \;\lor\;
                \bigl(a_k = (v_{k+1}, v_k) \;\land\; a_k \in E^{\sm D}\bigr)\,)
            \]
            with $v_0 \in A$ and $v_n \in J \cup B$ (the asymmetric right-endpoint inclusion of def \ref{def:sigma_separation}, item~1, instantiated on the CDMG $G^{\sm D}$ reads ``$v_n \in J^{\sm D} \cup B$''; using $J^{\sm D} = J$ recorded above, this is exactly $v_n \in J \cup B$), the walk $\pi$ is $C$-$\sigma$-blocked in $G^{\sm D}$ in the sense of def \ref{def:sigma_blocking} \emph{applied to the CDMG $G^{\sm D}$ in place of $G$} (``$C$-$\sigma$-blocked walk'' paragraph): there exists a position $k \in \{0, 1, \dots, n\}$ such that either ($k$ is a collider on $\pi$ in the sense of def \ref{def:collider_noncollider}, item~ii.\ computed relative to the edges of $G^{\sm D}$, and $v_k \notin \Anc^{G^{\sm D}}(C)$), or ($k$ is a blockable non-collider on $\pi$ in the sense of def \ref{def:unblockable_noncollider} computed relative to the edges of $G^{\sm D}$, and $v_k \in C$); the ancestor set on the right-hand side of each existential clause is $\Anc^{G^{\sm D}}(C) = \bigcup_{c \in C} \Anc^{G^{\sm D}}(c)$ from def \ref{def_3_5}, item~iv., \emph{computed in $G^{\sm D}$} (not in $G$).
    \end{enumerate}

    \emph{Reading of the equivalence.}
    Statements (1) and (2) are, respectively, the unfolded forms of ``$A$ is $i\sigma$-separated from $B$ given $C$ in $G$'' and ``$A$ is $i\sigma$-separated from $B$ given $C$ in $G^{\sm D}$'' under def \ref{def:sigma_separation}, item~1, instantiated on the two distinct CDMGs $G$ and $G^{\sm D}$ (the LN states the equivalence using the $\isPerp$ notation of def \ref{def:sigma_separation}, item~1, which is in this rewrite expanded inline as~(1) and~(2) so the spec is self-contained). The lemma asserts that these two universal statements are equivalent, despite the fact that switching from $G$ to $G^{\sm D}$ in general changes the carrier (from $J \cup V$ to $J \cup (V \sm D)$), the edge sets (from $E, L$ to $E^{\sm D}, L^{\sm D}$), the class of walks quantified over, the collider / blockable-non-collider classifications of positions on those walks (each computed relative to the edges of the CDMG in question), and the ancestor set used in the collider clause of $\sigma$-blocking (from $\Anc^G(C)$ to $\Anc^{G^{\sm D}}(C)$); the only data shared verbatim across the two sides are the input-node set $J^{\sm D} = J$, the triple $(A, B, C)$, and the conditioning role of $C$ inside the $\sigma$-blocking predicate.
\end{Lem}

\end{document}
